\section{Inside the gap there is no lattice}\label{sec:pointcount}

Section~\ref{sec:rank} put a rank-$m$ lattice at every integer $m\ge2$ and the
family in the first gap $m\in(1,2]$. This section shows that the lattice does not
interpolate into the gap, and identifies exactly what stops it.

\subsection{The interpolation exists and is forced}\label{sec:interpolation}

Write $q=e^{-\pi t}$ and $\theta=\sum_{k\in\Z}q^{k^{2}}=1+u$ with
$u=2q+2q^{4}+2q^{9}+\dots$, and for $j\ge1$ let $R_j(n)$ be the number of
\emph{ordered} $j$-tuples of \emph{positive} squares summing to $n$, so that
$u^{j}=\sum_n2^{j}R_j(n)q^{n}$. Let
\begin{equation}\label{eq:smin}
 s(n)=\min\{j:R_j(n)>0\}
\end{equation}
be the least number of positive squares summing to $n$.

\begin{proposition}[The interpolation {\cite[Prop.~3.1]{RankNote}}]\label{prop:interp}
$\theta$ does not vanish on the upper half-plane, so $\theta^{m}=\exp(m\log\theta)$
is well defined there for every real $m$, and
\begin{equation}\label{eq:rmn}
 \theta^{m}=\sum_{j\ge0}\binom mj u^{j},\qquad
 r_m(n)=\sum_{j=1}^{n}\binom mj\,2^{j}\,R_j(n),
\end{equation}
the sum terminating because $R_j(n)=0$ for $j>n$. Each $r_m(n)$ is a polynomial
in $m$ of degree at most $n$ which at every integer $m\ge0$ is the representation
number of $\Z^{m}$; being of degree $\le n$ and matching at $m=0,\dots,n$, it is
the unique such interpolation.
\end{proposition}

\subsection{The sign law, the free region, and the forcing}\label{sec:signlaw}

\begin{proposition}[The sign law]\label{prop:signlaw}
Let $k\ge0$ be an integer and $m\in(k,k+1)$. Then
\begin{equation}\label{eq:signlaw}
 \sgn\binom mj=(-1)^{\max(0,\;j-k-1)},
\end{equation}
that is: $\binom mj>0$ for $1\le j\le k+1$, $\binom m{k+2}<0$, and the sign
alternates thereafter. No $\binom mj$ vanishes.
\end{proposition}

\begin{proof}
$\binom mj=\frac1{j!}\prod_{i=0}^{j-1}(m-i)$, and for integer $i$ one has
$m-i<0$ exactly when $i\ge k+1$; among $i=0,\dots,j-1$ the number of such $i$ is
$\max(0,j-1-k)$. None vanishes because $m$ is not an integer.
\end{proof}

\begin{proposition}[The free region]\label{prop:free}
For $m\in(k,k+1)$ and $1\le n\le k+1$, $r_m(n)>0$.
\end{proposition}

\begin{proof}
$R_j(n)=0$ for $j>n$, so every contributing $j$ satisfies $j\le n\le k+1$ and has
$\binom mj>0$ by Proposition~\ref{prop:signlaw}; and $R_n(n)=1$, so the sum is
not empty.
\end{proof}

So negativity can begin only at $n\ge k+2$, and the first place at which it is
\emph{forced} --- where no positive term is present at all --- is governed by an
arithmetic quantity.

\begin{proposition}[Where the forcing lives, and where it stops]\label{prop:Nk}
Put $N_k=\min\{n\ge1:s(n)\ge k+2\}$. Then
\begin{equation}\label{eq:Nk}
 N_0=2,\qquad N_1=3,\qquad N_2=7,
\end{equation}
with $s(N_k)=k+2$ in each case; and $N_k$ \emph{does not exist} for any $k\ge3$,
because Lagrange's four-square theorem \cite[Thm.~369]{HardyWright} gives
$s(n)\le4$ for every $n\ge1$.
\end{proposition}

\begin{proof}
The three values are a finite verification. $2$ is not a square and $2=1+1$, so
$s(2)=2$. $3$ is not a square, and $3-1=2$ is not a square, so $3$ is not a sum
of two positive squares; $3=1+1+1$, so $s(3)=3$. $7$ is not a square; $7-1=6$ and
$7-4=3$ are not squares, so $7$ is not a sum of two positive squares;
$7-1-1=5$, $7-1-4=2$ and $7-4-4<0$ are not squares, so $7$ is not a sum of three
positive squares; and $7=4+1+1+1$, so $s(7)=4$. For the last clause, Lagrange
writes every $n$ as a sum of four squares; discarding the zero entries leaves at
most four positive ones.
\end{proof}

\begin{proposition}[The forced coefficient]\label{prop:forced}
For $k=0,1,2$ and every $m\in(k,k+1)$,
\begin{equation}\label{eq:forced}
 r_m(N_k)<0,\qquad\text{while}\qquad r_m(n)>0\ \text{ for every }1\le n<N_k .
\end{equation}
Explicitly $r_m(2)=2m(m-1)$, $r_m(3)=\frac43m(m-1)(m-2)$ and
$r_m(7)=64\binom m4+128\binom m7$.
\end{proposition}

\begin{proof}
At $n=N_k$ every contributing $j$ has $R_j(N_k)>0$, hence $j\ge s(N_k)=k+2$, so
every term carries a binomial of the signs of
Proposition~\ref{prop:signlaw} at or beyond $j=k+2$.

For $k=0$ and $k=1$ only one $j$ contributes at all: $R_j(2)=0$ except
$R_2(2)=1$, and $R_j(3)=0$ except $R_3(3)=1$. Hence $r_m(2)=4\binom m2$ and
$r_m(3)=8\binom m3$, both negative on their gap.

For $k=2$, $R_j(7)=0$ except $R_4(7)=4$ and $R_7(7)=1$, so
$r_m(7)=64\binom m4+128\binom m7$ and
\[
 r_m(7)=m(m-1)(m-2)(m-3)\Big[\tfrac83+\tfrac8{315}(m-4)(m-5)(m-6)\Big].
\]
On $(2,3)$ the prefactor is negative, and $(m-4)(m-5)(m-6)\in(-24,-6)$, so the
bracket lies in $\big(\frac83-\frac{192}{315},\ \frac83-\frac{48}{315}\big)
\subset(2.05,2.52)$ and is positive; hence $r_m(7)<0$.

For the second clause, $n\le k+1$ is Proposition~\ref{prop:free}, which settles
$n=1$ for $k=0$, $n\le2$ for $k=1$ and $n\le3$ for $k=2$. Three cases remain, all
on $(2,3)$. From $R_1(4)=R_4(4)=1$,
$r_m(4)=2m+16\binom m4=2m\big[1+\frac13(m-1)(m-2)(m-3)\big]$, and
$(m-1)(m-2)(m-3)\in(-2,0)$ there, so the bracket exceeds $\frac13$. From
$R_2(5)=2$ and $R_5(5)=1$, $r_m(5)=8\binom m2+32\binom m5$, and on $(2,3)$ both
binomials are positive by Proposition~\ref{prop:signlaw}. From $R_3(6)=3$ and
$R_6(6)=1$, $r_m(6)=24\binom m3+64\binom m6$ with $\binom m3>0$ and
$\binom m6<0$, and
$\frac{64\binom m6}{24\binom m3}=\frac83\cdot\frac{(m-3)(m-4)(m-5)}{120}$, whose
modulus is below $\frac83\cdot\frac6{120}=\frac2{15}<1$.
\end{proof}

\begin{corollary}[No lattice in the gap]\label{cor:nolattice}
For every $m\in(1,2)$ --- that is, every $d$ strictly inside the family's range
--- $\theta(it)^{m}=\sum_nr_m(n)e^{-\pi nt}$ is \emph{not} completely monotone:
being a Laplace expansion with discrete support, its representing measure is
forced by uniqueness to be $\sum_nr_m(n)\delta_{\pi n}$, and $r_m(3)<0$. There is
no positive measure and hence no point set.
\end{corollary}

\begin{remark}[What the mechanism is, and what it is not]\label{rem:lagrange}
Propositions~\ref{prop:signlaw}--\ref{prop:forced} say that the point count's
failure is a statement about \emph{sums of squares}: on the gap $(k,k+1)$ it is
forced by the least integer needing $k+2$ positive squares, and it stops being
forced at $k=3$ for one reason only, Lagrange's theorem. The scan reported in
\cite{RankNote} found the first negative coefficient at $n=2$ on $(0,1)$, $n=3$
on $(1,2)$, $n=7$ on $(2,3)$, at varying $n$ on $(3,4)$, and none at all for
$m\ge4$; the first three of those are now Proposition~\ref{prop:forced} and the
last is Proposition~\ref{prop:Nk}.

What is \emph{not} covered is negativity produced by competition rather than by
forcing. On $(3,4)$ nothing is forced, yet negative coefficients occur: there the
$j=k+2$ terms compete with positive ones instead of standing alone. Two
statements are therefore recorded as conjectures and are not used anywhere in
this manuscript: that $r_m(n)\ge0$ for every $n$ and every real $m\ge4$, and that
the negative coefficients on $(3,4)$ arise from competition. Both are questions
about sums of squares rather than about $\zeta$.
\end{remark}
